%! AP Statistics — Unit 3, Lesson 3.1 — Homework
\documentclass[10pt]{article}
\usepackage{apstats-article}
\usepackage{apstats-boxes}

\begin{document}

\pageheader{Unit 3, Lesson 3.1}{Homework}
\namedateperiod

\vspace{0.1in}

\begin{practicebox}
\small
For each survey below, (a) classify it as trustworthy or not trustworthy,
(b) identify whose opinion is over- or under-represented, and
(c) explain how the results might be distorted.

\vspace{8pt}
\textbf{1.} A magazine publishes an online poll asking readers whether they support
a new environmental regulation. Over 10{,}000 people vote.

\begin{enumerate}[label=(\alph*), leftmargin=2em, itemsep=8pt]
  \item Trustworthy / Not Trustworthy \quad \writeline
  \item Over/under-represented: \writeline
  \item Distortion: \writeline\\[1pt]\writeline
\end{enumerate}

\vspace{8pt}
\textbf{2.} A health researcher randomly selects 150 adults from a hospital's patient
database and asks about their exercise habits.

\begin{enumerate}[label=(\alph*), leftmargin=2em, itemsep=8pt]
  \item Trustworthy / Not Trustworthy \quad \writeline
  \item Who is included, and is there any concern? \writeline\\[1pt]\writeline
\end{enumerate}

\vspace{8pt}
\textbf{3.} A company surveys its own employees to determine what the average worker
in the city earns.

\begin{enumerate}[label=(\alph*), leftmargin=2em, itemsep=8pt]
  \item Trustworthy / Not Trustworthy \quad \writeline
  \item Whose opinion is missing? \writeline
  \item How might the result be distorted? \writeline\\[1pt]\writeline
\end{enumerate}

\vspace{8pt}
\textbf{4.} A researcher mails surveys to a random sample of 2{,}000 registered
voters. Only 300 respond.

\begin{enumerate}[label=(\alph*), leftmargin=2em, itemsep=8pt]
  \item Was the initial selection trustworthy? \blank{3cm}
  \item What concern arises from the low response rate? \writeline\\[1pt]\writeline
  \item Is a bigger (non-random) sample better? Explain. \writeline\\[1pt]\writeline
\end{enumerate}
\end{practicebox}

\end{document}
